\setlength{\parindent}{0em}
\begin{abstract}
A novel cross-layer optimization (\myCLO{}) scheme to
significantly enhance the spectral efficiency (\mySE{}) of a secondary
user (\mySU{}) transmission in an underlay cognitive radio system
(\myUCRS{}) is proposed. The scheme jointly optimizes the physical
layer (\myPHY{}) and data link control layer (\myDLC{}) of the \mySU{}~trans-\linebreak
mission. At the \myPHY{}, a joint \!bit and power loading \!(\myJBPL{})~with
coding adaptation (\myCA{}) using rate-matched low-density-parity-
check (\myLDPC) coding is employed. A truncated hybrid automatic
repeat request (\myTHARQ{}) using chase combining (\myCC{}), known~as
\myARQ{} type \myII{}, is implemented at the \myDLC{}.An analytic lower~bound\linebreak
on \!the \!\mySE{} for coded transmission is derived and maximized subject
to different quality-of-service (\myQoS{}) and stochastic chance-based
constraints as well as third-party colored Gaussian interference
at the \mySU{} receiver. The system performance is compared in terms
of the number of required retransmissions and the achievable \mySU{}
\mySE{} for different code rates. A significant improvement in system
performance is observed by using \myCC{} for the \mySU{} transmission.
\end{abstract}

